% Generated draft: requires mathematical and source review.
\documentclass[11pt]{article}
\usepackage[margin=25mm]{geometry}
\usepackage{fontspec}
\setmainfont{Latin Modern Roman}
\usepackage{amsmath,amssymb,mathtools,unicode-math}
\setmathfont{Latin Modern Math}
\usepackage{enumitem,needspace,xurl,hyperref}
\setlength{\parindent}{0pt}
\setlength{\parskip}{5pt}
\setlength{\emergencystretch}{3em}
\newcommand{\N}{\mathbb{N}}
\newcommand{\Z}{\mathbb{Z}}
\newcommand{\Q}{\mathbb{Q}}
\newcommand{\R}{\mathbb{R}}
\newcommand{\C}{\mathbb{C}}
\newcommand{\F}{\mathbb{F}}
\newcommand{\A}{\mathbb{A}}
\newcommand{\PP}{\mathbb{P}}
\newcommand{\E}{\mathbb{E}}
\newcommand{\eps}{\varepsilon}
\newcommand{\dd}{\,\mathrm{d}}
\begin{document}
\section*{UnsolvedMath ID: 701188}
\textbf{Source problem code:} catalogue entry
\subsection{Definitions and mathematical statement}
A group \(G\) is \emph{periodic} if every element of \(G\) has finite order. It is \emph{locally finite} if every subgroup generated by finitely many elements is finite. For a subgroup \(H\leq G\), its centralizer is
\[
C_G(H)=\{g\in G: gh=hg\text{ for every }h\in H\}.
\]
A noncyclic subgroup of order \(4\) is isomorphic to \((\Z/2\Z)\times(\Z/2\Z)\).

\textbf{Question.} For every periodic group \(G\), does the existence of a noncyclic subgroup \(H\leq G\) satisfying
\[
|H|=4\qquad\text{and}\qquad C_G(H)=H
\]
imply that \(G\) is locally finite?
\subsection{Short English statement}
Must every periodic group containing a noncyclic subgroup of order \(4\) whose centralizer is exactly that subgroup be locally finite?
\subsection{Sources}
\begin{itemize}
\item Local UnsolvedMath registry: permanent ID \(701188\), source problem code \texttt{catalogue entry}, title \emph{Kourovka Notebook Problem 15.100}.
\item Supplied attribution: A. K. Shl\"epkin, proposed in 2002.
\item Supplied reference: \emph{The Kourovka Notebook (Unsolved Problems in Group Theory)}, problem 15.100, 15th issue (2002); also identified in the record with the 21st issue edition (2026), \texttt{arXiv:1401.0300v48}. The collection description names E. I. Khukhro and V. D. Mazurov as editors.
\end{itemize}
\end{document}
